% !TEX program = LuaLaTeX+se
\documentclass[10pt,a6paper]{scrbook}
\pagestyle{empty}

% 1. MARGÓK BEÁLLÍTÁSA (Kisebb margó, hogy több férjen el)
\usepackage{geometry}
\geometry{
  top=1cm,
  bottom=1cm,
  left=1.5cm,
  right=1.5cm,
}

% betűtípus
\usepackage{fontspec}

\setmainfont{Linux Libertine O} % Linux Libertine betűtípus használata

% Nyelv

\usepackage{polyglossia}
\setdefaultlanguage{latin} % Alapértelmezett nyelv latin
\setotherlanguage{magyar} % Második nyelv magyar

% 3. Gregoriotex CSAK EZUTÁN
\usepackage{gregoriotex} 

% Opcionális: Szebb sorkizárás (segít eltüntetni a nagy lyukakat a szövegben)
\usepackage{microtype}

%\gredefsizedsymbol{VBarGoth}{greextra}{VWithBarGoth}
%\gredefsizedsymbol{RBarGoth}{greextra}{RWithBarGoth}

%\newcommand{\Vslash}{\VBarGoth{}.}
%\newcommand{\Rslash}{\RBarGoth{}.} 

\setkomafont{section}{\normalfont\centering\LARGE\scshape} 
\setcounter{secnumdepth}{-\maxdimen} 

% --- GREGORIO BEÁLLÍTÁSOK ---
% Alapméret
\grechangestaffsize{14}
% Iniciálé
\grechangedim{beforeinitialshift}{1.5mm}{scalable}
\grechangedim{afterinitialshift}{1.5mm}{scalable}
\grechangestyle{initial}{\fontsize{36}{36}\selectfont}
% Kottaszín
\gresetlinecolor{black}
% Kommentek ki/be (jobb felső sarok felirata)
\gresetheadercapture{commentary}{grecommentary}{string}
% Annotációk (pl. modus, tónus jelzés) stílusa
\grechangestyle{annotation}{\small\bfseries}
% --- TÉRKÖZÖK (A frissített nevekkel) ---
\grechangedim{baselineskip}{42pt}{scalable} % Kottasorok közötti függőleges távolság
\grechangedim{spaceabovelines}{1.5mm}{scalable} % Kotta FELETTI extra térköz
% JAVÍTVA: A modern Gregorio-ban a kotta alatti térköz a spacebeneathtext
\grechangedim{spacebeneathtext}{1.5mm}{scalable} 
% JAVÍTVA: A vonalvastagságot már nem a \grechangedim kezeli, hanem egy dedikált makró (alapértéke 17):
% \grechangestafflinethickness{17} 
\grechangedim{interglyphspace}{0.9mm}{scalable} % Hangjegyek közötti vízszintes távolság

%\addto\extrasmagyar{%
%  \hyphenpenalty=10000
%  \exhyphenpenalty=10000
%  \lefthyphenmin=62
%  \righthyphenmin=62
%  \emergencystretch=3em
%}

\begin{document}

\section{Ad Processionem In Fine Vesperarum}

\gregorioscore[a]{scores/antiphona/an--salve_regina--dominican}
\vspace{1em}
\noindent
\Vbar. Dignare me laudare te virgo sacráta - (T.P. Alleluia)\\
\textbf{\Rbar. Da mihi virtútem contra hostes tuos - (T.P. Alleluia)}
\vspace{1em}

\textbf{Orémus}

Concéde nos fámulos tuos, quǽsumus, Dómine Deus, perpétua mentis et córporis salúte gaudére: et gloriósa beatæ Maríæ semper Vírginis intercessióne a præsénti liberári tristía et ætérna pérfrui lætítia: Per Christum Dóminum nostrum \\
\textbf{\Rbar. Amen}
\vspace{1em}

\gregorioscore[a]{scores/antiphona/an--o_lumen--dominican}
\vspace{1em}
\noindent
\Vbar. Ora pro nobis, beáte Pater Dominíce - (T.P. Alleluia)\\
\textbf{\Rbar. Ut digni efficiámur promissiónibus Christi -}\\
{\raggedright (T.P. Alleluia.)\par}

\newpage
\textbf{Orémus}

Concéde quǽsumus omnípotens Deus: ut, qui peccatórum nostrórum póndere prémimur - beáti Domínici Confessóris tui, Patris nostri, patrocínio sublevémur: Per Christum Dóminum nostrum \\
\textbf{\Rbar. Amen}

\vspace{1em}
\noindent
\Vbar. Fidelium animae per nostrum. Dei requiescant in pace.\\
\textbf{\Rbar. Amen}

\section{Ima a hivatásokért}

\begin{magyar}
{\large
\hyphenpenalty=10000
\indent Urunk, Jézus Krisztus, Te egykor meghívtad Egyházadba Szent Domonkos családját, hogy hirdesse az evangéliumot.\\
\indent Kérünk, küldj ugyanúgy ma is apostolokat aratásodba. Adj nekik bátorságot, bölcsességet és kegyelmet, hogy az emberek előtt méltóképpen tegyenek tanúságot halálodról, feltámadásodról és dicsőséges eljöveteledről.\\
\indent Engedd, hogy legszentebb Édesanyádnak, Rendünk pártfogójának közbenjárására legyünk a hit védőivé, a közömbösséget és a bűn sötétségét szétoszlató fénnyé: te, aki mindig velünk vagy. Amen  
}
\end{magyar}

\end{document}
